% Pursuing Stacks (cotes 134-2 to 134-8): archive cote 134-6, site page 44 (PDF page 45), typescript page 362.
% typescript. Draft compared: ps5.tex:3530-3604.
% AI-transcribed by Claude Opus 5.5 (claude-opus-5-5, adaptive thinking, effort high), Message Batches API, batch ps3 (msgbatch_01XdFmuVVBAcgc7EFxdtCqsj), 2026-11-12; reviewed against both supplied scan halves; see CHECK_DECISIONS.tsv.
\scanpage{134-6/44}
\mspage{362}

of $B^{\hand{\wedge}}_{ab}$ . I will first describe this property or \inserted[typed]{possible} function of a chain
complex in Add(B), forgetting for the time being the category $A=B^{o}$
and the homotopy types defined by objects X in $A^{\hand{\wedge}}$. Onc\over[typed]{e}{e} this property
is \over[typed]{o}{w}ell understood, it will be time to show it implies the previous
one relative to A and objects of $A^{\hand{\wedge}}$, and presumably is even equivalent
with it.

First we'll have to interprete the category Add(B), which was
constructed somewhat ``abstractly'' in section 93 (as the solution of a
universal problem stated in section 92), as a full subcategory of
the category $B^{\hand{\wedge}}_{ab}$ of abelian presheaves on B. It will be useful to
keep in mind the following diagram of canonical functors

% Hand-drawn arrows are marked blue; typed objects remain plain.
(2) \diagram{\begin{tikzcd}
\ensuremath{B} \arrow[color=blue,r, "\hand{\alpha_B}"] \arrow[color=blue,d] \arrow[color=blue,dr, "\hand{\beta_B}"] \& \ensuremath{B^{\hand{\wedge}}} \arrow[color=blue,d, "Wh_B"] \\
\ensuremath{Add(B)} \arrow[color=blue,r, "\hand{\gamma_B}"'] \& \ensuremath{B^{\hand{\wedge}}_{ab}}
\end{tikzcd}} ,

where $\hand{\alpha}_B$ is the canonical inclusion, $Wh_B$ is the abelianization func-%
tor, $\hand{\beta}_B$ the composition of the two, and $\hand{\gamma}_B$ the \uline{additive} functor
factoring $\hand{\beta}_B$ , in virtue of the universal property of Add(B). This
functor is defined up to canonical isomorphism, and the lower triangle
of (2) is commutative, up to a given commutativity isomorphism. Also,
we'll use the composition of the following sequence of canonical equi-%
valences of categories:
\[ B^{\hand{\wedge}}_{ab} = \underline{Hom}(B^{o},(Ab)) \hand{\xrightarrow{\approx}} \underline{Hom}(B,(Ab)^{o})^{o} \hand{\xrightarrow{\approx}} \underline{Homadd}(Add(B),(Ab)^{o})^{\hand{o}}\deleted{\illegible{1}} \]
\[ \hand{\xrightarrow{\approx}} \underline{Homadd}(Add(B)^{o},(Ab)) , \]
i.e. a canonical equivalence of categories
\[ (3) \qquad B^{\hand{\wedge}}_{ab} \hand{\xrightarrow{\approx}} \underline{Homadd}(Add(B)^{o}, (Ab)) \quad , \quad F \hand{\mapsto} \overset{\hand{\sim}}{F} , \]
which \deleted{can be generalized into} \inserted{is a particular case} \hand{of}
\[ (3') \qquad \underline{Hom}(B^{o}, M) \hand{\xrightarrow{\approx}} \underline{Hom}(Add(B)^{o}, M ) \quad , \]
where M is any \xout{abelian} \inserted[typed]{additive} category. If F is an abelian presheaf on B
i.e. an object in the \xout{first} \inserted[typed]{left} hand side of (3), we'll denote by
\[ (4) \qquad \overset{\hand{\sim}}{F} : \ Add(B)^{o} \hand{\longrightarrow} (Ab) \ \deleted{\illegible{1}} \]
the corresponding additive functor. Now this functor can be interpre-%
ted very nicely in terms of the functor $\hand{\gamma}_B$ in (2), by the canonical
isomorphism of abelian groups
\[ (5) \qquad \overset{\hand{\sim}}{F}(L) \hand{\xrightarrow{\sim}} Hom_{B^{\hand{\wedge}}}(\hand{\gamma}_B(L),F)\hand{,} \]

\margin{44}

% Visual collation: both supplied image halves inspected; see CHECK_DECISIONS_134_7_8.tsv.
